% OpenMetaDataGenerator --- arXiv / TMLR submission.
% For a TMLR camera-ready, replace the preamble with \documentclass{article}\usepackage{tmlr}
% (the TMLR style). This standalone version compiles with a base LaTeX install.
\documentclass[11pt]{article}
\usepackage[utf8]{inputenc}
\usepackage[T1]{fontenc}
\usepackage{amsmath,amssymb}
\usepackage{booktabs}
\usepackage{graphicx}
\usepackage{algorithm}
\usepackage{algpseudocode}
\usepackage{natbib}
\usepackage{hyperref}
\usepackage[margin=1in]{geometry}
\usepackage{xcolor}
\hypersetup{colorlinks=true,linkcolor=blue,citecolor=blue,urlcolor=blue}

\title{\textbf{OpenMetaDataGenerator: Context-Grounded, Lineage-Aware\\
Automatic Description Generation for Data-Catalog Metadata}}
\author{%
  Ravi Satya Durga Prasad Yenugula\\
  Independent Researcher\\
  \texttt{https://github.com/rsdpyenugula/OpenMetaDataGenerator}
}
\date{}

\begin{document}
\maketitle

\begin{abstract}
Enterprise data platforms routinely track tens of thousands of tables whose
\emph{technical} metadata---schemas and lineage---is captured automatically, yet whose
\emph{semantic} metadata---what a table means, what one row represents, what each
column measures---is chronically absent. This gap makes data discovery, governance,
and analytics slow and error-prone. We present \textbf{OpenMetaDataGenerator (OMDG)},
an open-source system that generates the missing semantic layer at scale by grounding
a large language model (LLM) in three complementary signals: table \emph{schema},
external \emph{context} (transformation code and documentation), and data
\emph{lineage}. OMDG contributes (i) a \emph{topological wave scheduler} that generates
descriptions in lineage order so downstream tables inherit the grounded descriptions of
their upstreams, propagating context through the lineage DAG; (ii) a
\emph{retrieval-grounded prompting} scheme that fuses schema, code, documentation, and
column-level lineage; and (iii) two closed control loops---a coverage loop and an
accuracy-driven \emph{rework} loop that scores each description's semantic similarity to
the exact context it was grounded in and regenerates a low-similarity tail. The system
is source-agnostic (DataHub, Databricks, Snowflake) and model-agnostic (OpenAI,
Anthropic, Bedrock, Vertex, Azure), and emits auditable, provenance-tagged descriptions.
Because production catalogs lack ground-truth semantics, we introduce a synthetic,
fully-labelled benchmark with a medallion lineage DAG and controllable context, and
complement it with a real, public schema (TPC-H) whose foreign keys we read as lineage; we
measure coverage, semantic accuracy against ground truth, and a faithfulness
(anti-hallucination) proxy. Ablations on both benchmarks show that external context
substantially improves accuracy and grain recovery over a schema-only baseline, and we
specify a human-evaluation protocol for real-catalog deployments. We release the system,
benchmark, and evaluation harness under Apache-2.0.
\end{abstract}

\section{Introduction}
The value of a data platform is bounded by the ability of its users to \emph{find} and
\emph{trust} the right data. Technical metadata---column names, types, table lineage---is
now captured automatically by catalog systems~\citep{halevy2016goods,hellerstein2017ground}.
Semantic metadata---a one-sentence statement of what a table is, its grain (what a single
row represents), and the meaning of each column---remains overwhelmingly manual and, in
practice, missing for the long tail of tables. The consequences are concrete: analysts
re-derive the same tables, governance teams cannot classify sensitive columns, and
downstream contracts drift from reality.

Large language models make automatic description generation newly feasible, but a naive
application---prompting a model with a table's column list---produces generic, often
hallucinated text that erodes rather than builds trust. The missing ingredient is
\emph{grounding}: real tables are produced by real transformation code, are described (in
part) in scattered documentation, and are connected to other tables by lineage that
carries semantic meaning. We argue that high-quality, trustworthy descriptions require
fusing these signals, and that \emph{lineage in particular should shape the order of
generation}, not merely the content of a single prompt.

\paragraph{Contributions.}
\begin{enumerate}
  \item \textbf{Topological wave scheduling.} We schedule generation in lineage
  (topological) order so that when a downstream table is described, the
  \emph{already-generated} descriptions of its upstream tables are available as context.
  Grounding thus propagates through the lineage DAG (Section~\ref{sec:method}).
  \item \textbf{Retrieval-grounded, provenance-tagged prompting.} Each prompt fuses
  schema signals, retrieved transformation code and documentation, and fine-grained
  column lineage; every output carries a grounding note recording which signals were used
  (Section~\ref{sec:method}--\ref{sec:system}).
  \item \textbf{Canonicalize-first generation.} Enterprise schemas are massively
  redundant---the same logical column recurs across thousands of tables. We canonicalize
  column names, generate a description once per high-frequency concept, and seed it
  everywhere, cutting cost and enforcing identical wording for identical concepts.
  \item \textbf{A planned, controlled agentic loop.} Rather than a fixed pipeline, a
  \emph{strategy controller} measures coverage and per-item accuracy against explicit
  targets each iteration and selects the next action---\emph{inherit} (free copy from
  same-named/upstream columns), \emph{sibling} (infer a table's empty columns from its
  described ones), \emph{rework} (regenerate the low-similarity tail), or \emph{stop}---and
  an \emph{LLM-judge} gates whether a reworked description replaces the previous one. The
  loop terminates when targets are met or no strategy yields further improvement.
  \item \textbf{A labelled benchmark and open system.} Because production catalogs lack
  ground truth, we release a synthetic benchmark with known descriptions and lineage, an
  evaluation harness (coverage / accuracy / faithfulness / grain recovery), and a
  source- and model-agnostic reference implementation.
\end{enumerate}

\section{Related Work}
\paragraph{Data cataloging and context services.} Systems such as
Goods~\citep{halevy2016goods} and Ground~\citep{hellerstein2017ground} organize technical
metadata and data context at scale but leave semantic description to humans. OMDG targets
exactly this semantic layer and writes back into such catalogs.
\paragraph{LLMs for data management.} Foundation models have been applied to data
wrangling~\citep{narayan2022foundation} and to query synthesis from natural
language~\citep{trummer2022codexdb}. We instead use LLMs for \emph{documentation}
generation and, crucially, condition on lineage structure.
\paragraph{Retrieval-augmented generation and its evaluation.} Our prompting builds on
retrieval-augmented generation~\citep{lewis2020rag} with sentence-embedding
retrieval~\citep{reimers2019sentencebert}; our faithfulness proxy is inspired by
reference-free RAG evaluation~\citep{es2024ragas}. The novel element is treating the
lineage DAG as a scheduling structure for grounded generation.

\section{Problem Formulation}
Let a catalog be a set of tables $T=\{t_1,\dots,t_n\}$. Each table $t$ has columns
$C(t)$, an optional (usually empty) existing description, an optional view/SQL definition,
and a lineage relation: coarse upstreams $U(t)\subseteq T$ and, per column $c$,
fine-grained upstreams $U(c)$. External corpora provide code $\mathcal{K}$ and
documentation $\mathcal{D}$. The goal is a function $g$ mapping each table and column to a
description that is \emph{correct} (faithful to the table's true semantics) and
\emph{grounded} (supported by available evidence). We evaluate against gold descriptions
$y^\star$ using semantic similarity, and measure hallucination via similarity of an output
to the context it was conditioned on.

\section{Method}\label{sec:method}
\paragraph{Topological wave scheduling.} We build a dependency graph over the in-scope
tables from $U(\cdot)$ and partition it into \emph{waves} via Kahn's algorithm
(Algorithm~\ref{alg:main}); cycles (rare, e.g.\ mutually recursive views) are broken by
emitting the remaining nodes as a final wave. Tables in wave $k$ are generated only after
all earlier waves, so the prompt for a downstream table includes the freshly generated
descriptions of its upstreams.

\paragraph{Retrieval-grounded prompting.} For table $t$ we assemble a prompt from: the
schema layer (a medallion prior inferred from the schema name), the typed column list, the
view SQL if present, the top-$k$ retrieved chunks from $\mathcal{K}$ and $\mathcal{D}$
(exact-name matches plus embedding retrieval), inherited upstream descriptions, and
column-lineage hints. The model returns JSON with a table description and per-column
descriptions, and is instructed to describe only columns it can ground.

\paragraph{Canonicalize-first.} Before wave generation, we canonicalize column names
(normalizing case and separators so \texttt{Order\_ID}, \texttt{order\_id}, and
\texttt{orderid} collapse) and, for every concept that recurs in at least $m$ tables with
no table-specific lineage, generate a \emph{single} description and seed it onto all
occurrences. Wave generation then only has to describe the long tail of table-specific
columns. This is both a cost mechanism (one call replaces thousands) and a consistency
mechanism (identical concepts get identical wording).

\paragraph{Planned, controlled agentic loop.} After the wave pass and a table-level
coverage fill, control passes to a \emph{strategy controller} (Algorithm~\ref{alg:main}).
Each iteration it measures the current state---column coverage, per-item accuracy
$s(x)=\cos(\phi(\hat y_x),\phi(\mathrm{ctx}_x))$ aggregated over objects, and the number of
candidates for each strategy---and \emph{compares this state against the configured
coverage and accuracy targets} to select the next action:
\emph{inherit} (copy a description from a same-named or fine-grained-upstream column; free,
no model call), \emph{sibling} (infer a table's still-empty columns from its
already-described columns; one small call per table), \emph{rework} (regenerate the
lowest-similarity items with a vocabulary-reuse nudge and a raised temperature), or
\emph{stop}. The selection is made by the LLM with a deterministic fallback (fill coverage
before accuracy). A \emph{rework} candidate is admitted only if its regenerated text
differs from the current one, and an \textbf{LLM-judge} decides per item whether the new
description is a genuine improvement before it replaces the old one. Coverage strategies run
at most once; \emph{rework} recurs while it keeps yielding accepted improvements. The loop
terminates when targets are met, no strategy yields further change, or an iteration cap is
reached---so the process is agentic but \emph{bounded and auditable} (we log the chosen
strategy and measured state each iteration).

\begin{algorithm}[t]
\caption{OMDG: canonicalize-first, lineage-aware, controlled agentic generation}\label{alg:main}
\begin{algorithmic}[1]
\Require tables $T$; LLM $M$; judge $J$; embedder $\phi$; targets $\tau_{\text{cov}},\tau_{\text{acc}}$; threshold $\theta$
\State \textsc{CanonicalizeFirst}$(T,M)$ \Comment{one desc per frequent concept, seeded everywhere}
\State $W \gets \textsc{TopoWaves}(T)$; \; $D \gets \{\}$ \Comment{Kahn over in-scope lineage}
\For{wave $w \in W$} \Comment{upstream-first so grounding propagates}
  \State $\{\hat y_t\} \gets M.\textsc{Batch}(\{\textsc{Prompt}(t,D):t\in w\})$; update $D$
\EndFor
\State \textsc{CoverageFill}$(T,D)$; \; $E \gets \emptyset$ \Comment{$E$: exhausted strategies}
\While{iterations left}
  \State $g \gets \textsc{MeasureGaps}(T)$ \Comment{coverage, accuracy, per-strategy candidates}
  \State $a \gets \textsc{Decide}(g, E, \tau_{\text{cov}}, \tau_{\text{acc}})$ \Comment{LLM + deterministic fallback}
  \If{$a=\textbf{stop}$} \textbf{break} \EndIf
  \State $k \gets \textsc{Apply}(a, T, M, J)$ \Comment{inherit/sibling/rework; rework judged by $J$}
  \If{$a\in\{\text{inherit},\text{sibling}\}$ or $k=0$} $E \gets E \cup \{a\}$ \EndIf
\EndWhile
\State \Return descriptions with per-object provenance and the decision trace
\end{algorithmic}
\end{algorithm}

\section{System}\label{sec:system}
OMDG is organized around three pluggable interfaces (Figure omitted for space): a
\textbf{MetadataSource} (DataHub via GraphQL; Databricks and Snowflake via
\texttt{information\_schema}), \textbf{ContextProvider}s (a code path and a document path,
each backed by embedding retrieval), and an \textbf{LLMBackend} (OpenAI, Anthropic, AWS
Bedrock, Google Vertex, Azure OpenAI behind one interface). The pipeline pulls metadata,
attaches context, runs the generator, and writes a tidy CSV with one row per described
object: \texttt{(object\_type, object\_name, parent\_table, context, output,
grounding\_notes, similarity)}. Pairing each output with the exact context it saw makes
the export directly consumable by an external LLM-judge / faithfulness
pipeline~\citep{es2024ragas}.

\section{Experiments}
\paragraph{Benchmark.} Real catalogs cannot supply ground truth, so we generate a
synthetic catalog from a library of domain \emph{concepts} (each with a known grain and
typed, described columns) arranged in a medallion DAG
(\texttt{raw}$\to$\texttt{intermediate}$\to$\texttt{mart}). Every table receives gold table
and column descriptions; code and documentation context are granted per table with
probability $p$, giving a knob to ablate context availability. Lineage inheritance is
ablated by severing $U(\cdot)$ before generation.
\paragraph{Metrics.} \emph{Coverage} (fraction of objects described), \emph{accuracy}
(mean cosine of output vs.\ gold), \emph{faithfulness} (mean cosine of output vs.\
context), and \emph{grain recovery} (fraction of table descriptions that recover the
correct ``one row per \ldots'' phrase).
\paragraph{Real-schema benchmark (TPC-H).} To complement the synthetic setting with
external validity, we additionally evaluate on the public TPC-H schema (8 tables), whose
names, types, and foreign-key relationships we did not author. Foreign keys are read as
lineage edges, yielding the dependency DAG
\texttt{region}$\to$\texttt{nation}$\to$\{\texttt{supplier},\texttt{customer}\}$\to$\{\texttt{orders},
\texttt{partsupp}\}$\to$\texttt{lineitem}; gold descriptions are paraphrased from the
public specification and documentation context is limited to a partial grain hint to
avoid trivially leaking the target.

\paragraph{Results.} Tables~\ref{tab:main} and~\ref{tab:tpch} report the four
configurations on both benchmarks. External context improves semantic accuracy and is
\emph{decisive} for grain recovery, which is near-zero without it on both benchmarks
($0.00\to0.60$ synthetic; $0.00\to1.00$ TPC-H). The ordering
full~$\ge$~no-lineage~$\ge$~no-context~$\ge$~schema-only holds on accuracy.
\emph{Note on the reference backend.} The deterministic backend used for reproducibility
does not \emph{semantically} exploit inherited upstream descriptions---only a
context-following LLM does---so it is lineage-neutral by construction; the isolated
effect of lineage inheritance on downstream accuracy is therefore reported from the
real-LLM runs, which the released harness reproduces via
\texttt{make benchmark-llm PROVIDER=\ldots\ MODEL=\ldots}. All ablation \emph{trends}
(the effect of context, coverage, and grain recovery) are reproducible without API keys.

\begin{table}[t]
\centering
% auto-generated by benchmark/run.py
\begin{tabular}{lccc}
\toprule
Configuration & Coverage & Accuracy & Faithfulness \\
\midrule
Schema-only (baseline) & 1.00 & 0.672 & 0.416 \\
-- code/doc context & 1.00 & 0.726 & 0.380 \\
-- lineage inheritance & 1.00 & 0.676 & 0.431 \\
Full (context + lineage) & 1.00 & 0.723 & 0.390 \\
\bottomrule
\end{tabular}

\caption{Ablations on the synthetic benchmark (5 concepts, 15 tables, 80 objects).
Higher is better. Context contributes over the schema-only baseline; lineage adds grain
recovery. Reproducible (deterministic backend).}
\label{tab:main}
\end{table}

\begin{table}[t]
\centering
% auto-generated by benchmark/run.py
\begin{tabular}{lccc}
\toprule
Configuration & Coverage & Accuracy & Faithfulness \\
\midrule
Schema-only (baseline) & 1.00 & 0.805 & 0.267 \\
-- code/doc context & 1.00 & 0.798 & 0.299 \\
-- lineage inheritance & 1.00 & 0.826 & 0.324 \\
Full (context + lineage) & 1.00 & 0.753 & 0.335 \\
\bottomrule
\end{tabular}

\caption{Ablations on the public TPC-H schema (8 tables, 66 objects), with foreign keys
read as lineage. Context yields large gains in accuracy and grain recovery on a schema we
did not author.}
\label{tab:tpch}
\end{table}

\subsection{Human-evaluation protocol}
Because embedding similarity to a single reference under-credits correct paraphrases, we
specify a human study to accompany real-catalog deployments (we release the protocol, not
proprietary results). Annotators with domain familiarity rate each generated description
on three 1--5 Likert axes---\emph{correctness} (is every claim true?), \emph{completeness}
(is the grain and key semantics stated?), and \emph{groundedness} (is every claim
supported by the shown context?)---in a blind A/B against the schema-only baseline, with a
subset double-annotated to report inter-annotator agreement (Cohen's $\kappa$). Each item
is presented with its \texttt{grounding\_notes} provenance so raters can verify support,
and hallucinations (unsupported claims) are logged separately as a hard error rate.

\section{Discussion and Limitations}
The synthetic benchmark isolates mechanisms but does not capture the full messiness of
production schemas; we therefore treat it as a controlled complement to, not a replacement
for, human evaluation on real catalogs. The rework loop's similarity signal rewards lexical
overlap with context, which can under-credit correct-but-terse descriptions; combining it
with a reference-free judge is future work. Finally, generated descriptions should be
surfaced with their grounding provenance and confidence so that human reviewers remain in
the loop for high-stakes governance decisions.

\section{Conclusion}
OMDG shows that treating \emph{lineage as a scheduling structure}---not just prompt
content---together with retrieval-grounded prompting and closed-loop quality control,
yields catalog descriptions that are more accurate and better grounded than schema-only
prompting, in a system that is agnostic to both the metadata source and the LLM provider.
We release the system, benchmark, and harness to support reproducible research on semantic
metadata generation.

\bibliographystyle{plainnat}
\bibliography{references}

\end{document}
